\section{One quadratic form detects both critical-line conjectures}
\label{sec:prime}
Choose $Q$ to be either $Q_1$ or $Q_2$, once and for all. Write
\[
 S_Q(x)=\sum_{\substack{p\le x,\ p\ne47\\ p=Q(u,v)\text{ for some }u,v\in\Z}}
   (1-p/x)\log p .
\]
The sum counts each prime once. The appropriate conjugacy class $C$ comprises two of the ten group elements, so the unsmoothed Chebyshev density is $1/5$. The linear smoothing halves the main term.

\begin{theorem}\label{th:primes}
Use the verified Dirichlet GRH for $\zeta$ and $L(s,\chi_{-47})$. For either single choice of $Q$ the following are equivalent.
\begin{enumerate}
\item The nontrivial zeros of both $L_1$ and $L_2$ lie on $\Re s=1/2$.
\item For every $\epsilon>0$ there is $A_\epsilon$ such that, eventually,
\[
                 S_Q(x)\le x/10+A_\epsilon x^{1/2+\epsilon}.
\]
\item The same statement holds with the lower bound
$S_Q(x)\ge x/10-A_\epsilon x^{1/2+\epsilon}$.
\end{enumerate}
In particular either direction of error suffices by itself, and common zeros of the two factors cannot defeat the test.
\end{theorem}

\paragraph{Proof.}
Let $r$ generate the rotations of the dihedral group, labeled so the ideal class of $Q_1$ corresponds to $r$ up to inversion. By ideal class reciprocity, unramified rational primes represented by the form are exactly those with Frobenius in $C=\{r,r^{-1}\}$, or $C=\{r^2,r^{-2}\}$ for the second form. Put $a=(\sqrt5-1)/2$ and $b=-(\sqrt5+1)/2$. For the first class, character column orthogonality gives
\[
 5\mathbf{1}_C=\mathbf 1+\mathrm{sgn}+a\chi_1+b\chi_2 .
\]
Here $\chi_j$ is the Artin character of $f_j$. For the second class interchange $a,b$. In particular $a+b=-1$, and $a,b$ are linearly independent over $\Q$.

Let $U(x)$ sum $\mathbf 1_C(g_p^k)(1-p^k/x)\log p$ over $p^k\le x$, $p\ne47$. Chebyshev's estimate bounds the difference from $S_Q$ by $O(\sqrt x+x^{1/3}\log x)=O(\sqrt x)$. For $\Re s>1$,
\begin{align*}
 \int_1^\infty U(x)x^{-s-1}\,\dd x
 &=\frac{D(s)}{s(s+1)},\\
 5D(s)&=-\frac{\zeta'}\zeta(s)-\frac{L'}L(s,\chi_{-47})
         -a\frac{L'_1}{L_1}(s)-b\frac{L'_2}{L_2}(s)+B(s).
\end{align*}
The ramified correction $B$ is holomorphic on $\Re s>0$, since the finite-image local eigenvalues have absolute value one on inertia invariants. The residue of the integral at $s=1$ is $1/10$. Subtract $1/[10(s-1)]$ and call the resulting meromorphic function $M(s)$. It is regular at every real point of $(1/2,\infty)$ by the Dirichlet input, Theorem~\ref{th:real}, and the cancellation at one. At a candidate nonreal zero $\rho$ with $\Re\rho>1/2$, its residue equals
\[
     -\frac{a\,\operatorname{ord}_\rho L_1+b\,\operatorname{ord}_\rho L_2}
            {5\rho(\rho+1)} .
\]
The numerator can vanish only when both integer multiplicities vanish.

Take $\sigma_0=1/2+\epsilon<1$. Either one-sided estimate supplies a sign and a constant $B_0$ for which
$V(x)=B_0x^{\sigma_0}\pm (U(x)-x/10)$ is nonnegative on $[1,\infty)$, after enlargement on a compact interval. Its Mellin transform continues as $B_0/(s-\sigma_0)\pm M(s)$ and thus is regular on the real interval $(\sigma_0,\infty)$. For completeness, a nonnegative Mellin integrand with finite abscissa $\alpha$ has a singularity there. If regular there, choose a real $s_0>\alpha$ sufficiently close that the Taylor disk of the transform at $s_0$ reaches left of $\alpha$. Evaluating at $s_0-\delta$, the integrated Taylor terms are nonnegative powers of $\delta\log x$ times the integrand. Monotone convergence then gives convergence left of the abscissa, a contradiction. Thus $\alpha\le\sigma_0$ (or convergence holds everywhere). The transform is holomorphic in the open half-plane to its right. The displayed residue excludes all zeros there. Let $\epsilon\downarrow0$ and use the Fricke functional equations to reflect the result.

Conversely Artin factorization gives $\zeta_H=\zeta L(s,\chi_{-47})L_1^2L_2^2$. Under (1) and Dirichlet GRH, effective Chebotarev for this fixed field gives the ordinary weighted count $x/5+O(\sqrt x\log^2(2x))$. Partial integration with $1-p/x$ gives both one-sided estimates. See Greni\'e and Molteni, \href{https://arxiv.org/abs/1709.07609}{An effective Chebotarev density theorem under GRH}.\hfill$\square$

\paragraph{What the arithmetic construction must still do.}
Even in the mean-square sense no rational congruence weight reproduces this class. Let $M$ be any modulus and $v$ any function of a reduced residue class modulo $M$. In the compositum of $H$ and $\Q(\zeta_{\operatorname{lcm}(47,M)})$, the intersection of $H$ with the cyclotomic field is $K$: it contains $K$, and $K$ is the maximal abelian subextension of $H/\Q$. Chebotarev consequently gives a uniform distribution on the five rotations after fixing a compatible residue class that splits in $K$. The conditional probability of $C$ there is $2/5$, and it is zero when the rational prime is inert. The best constant approximation on each fiber is its conditional mean. Therefore
\[
 \lim_{X\to\infty}\frac1{\pi(X)}\sum_{p\le X}^{*}
  |\mathbf1_C(g_p)-v(p\bmod M)|^2\ \ge\
 \frac12\frac25\frac35=\frac3{25}.
\]
The star omits the finitely many excluded primes. Equality is achieved for $M=47$ using that conditional mean. This obstruction concerns substitution of residue weights. It still allows a genuinely new theta-lift construction or cancellations between different arithmetic estimates. For example, an exhaustive enumeration through $10^6$ finds $7{,}819$, $15{,}692$ and $15{,}690$ primes represented by $Q_0,Q_1,Q_2$, respectively, with $47$ omitted. The three disjoint sets exhaust exactly the primes with quadratic symbol $+1$. For the last two classes the centered values $(S_Q(x)-x/10)/\sqrt x$ are about $-0.04823$ and $-0.11582$ at this cutoff. The counts check the identification of the classes and the sample illustrates the normalization. The theorem requires a bound at all sufficiently large $x$.
